\subsection{Proper partial maps}

In this section, X and Y are locally compact topological spaces. We denote by \(X'\) (resp. \(Y'\) ) the compact space obtained from X (resp. Y) by adjoining a point at infinity \(\omega_{X}\) (resp. \(\omega_{Y}\) ) (TG, I, p.~67--68). We identify \(X'\) and \(Y'\) with \(\mathsf{X}'(\mathcal{C}_{0}(\mathsf{X}))\) and \(\mathsf{X}'(\mathcal{C}_{0}(\mathsf{Y}))\) respectively (Corollary 2 of I, p.~32).

DEFINITION 1. --- A proper partial map from X to Y is a correspondence \(f = (\Gamma, \mathrm{X}, \mathrm{Y})\) (E, II, p.~10, def. 2) between X and Y such that

\begin{enumerate}
\def\labelenumi{(\roman{enumi})}
\item
  The graph \(\Gamma\) is functional;
\item
  The domain of definition of f is an open set U of X;
\item
  The map \(x\mapsto f(x)\) from U to Y is proper.
\end{enumerate}

The identity map on X is a proper partial map from X into X. Let Z be a locally compact topological space and \(f\) (resp. \(g\) ) a proper partial map from X to Y (resp. from Y to Z). Then the composite correspondence \(g \circ f\) (E, II, p.~11, def. 6) is a proper partial map from X into Z (TG, I, p.~72, prop. 3, and p.~73, prop. 5).

Lemma 1. --- For every proper partial map f from X to Y, with domain of definition U, denote by \(\widetilde{f}\) the mapping from X' to Y' defined by \(\widetilde{f}(x) = f(x)\) if \(x \in \mathrm{U}\) and \(\widetilde{f}(x) = \omega_{\mathrm{Y}}\) if \(x \notin \mathrm{U}\) ; it is continuous.

The map \(f \mapsto \widetilde{f}\) is a bijection between the set of proper partial maps \(f\) from X to Y and the set of continuous maps \(g\) from X' to Y' such that \(g(\omega_{\mathrm{X}}) = \omega_{\mathrm{Y}}\) .

Let \(f\) a proper partial map from X to Y, and let U be its domain. Let us prove that the map \(\widetilde{f}\) is continuous. It is continuous at every point of U, since U is open in \(\mathrm{X}'\) Let us prove that it is also continuous at every point \(x\) of \(\mathrm{X}' - \mathrm{U}\) ; we then have \(\widetilde{f}(x) = \omega_{\mathrm{Y}}\) Let V be an open neighborhood of \(\omega_{\mathrm{Y}}\) into \(\mathrm{Y}'\) ; let us prove that \(\widetilde{f}^{-1}(\mathrm{V})\) is a neighborhood of \(x\) By definition of the topological space \(\mathrm{Y}'\) , one may assume that V is of the form \(Y^{\prime}-K\) where K is a compact subset of Y. Since f defines a proper map from U to Y, the set \(f^{-1}(K)\) is compact in U (TG, I, p.~77, prop. 6), hence in \(X^{\prime}\) This is in particular a closed subset of \(X^{\prime}\) and \(\tilde{f}^{-1}(V)=X^{\prime}-f^{-1}(K)\) is an open subset of \(X^{\prime}\) and is therefore a neighborhood of x.

Conversely, let \(g \colon X' \to Y'\) be a continuous map such that \(g(\omega_{\mathrm{X}}) = \omega_{\mathrm{Y}}\) and \(\Gamma_g \subset X' \times Y'\) its graph. The set \(U = X - \overline{g}^{-1}(\omega_{\mathrm{Y}})\) is open in \(X\) . The correspondence \(f = (\Gamma_g \cap (U \times Y), X, Y)\) is a proper partial map from \(X\) into \(Y\) such that \(\widetilde{f} = g\) , and it is the only one.

We will identify the proper partial maps from X to Y with the continuous maps from \(X'\) into \(Y'\) that map \(\omega_{X}\) to \(\omega_{Y}\) . In particular, the proper maps from X to Y are the proper partial maps with domain X; they are identified with the continuous maps f from \(X'\) into \(Y'\) such that \(f^{-1}(\omega_{\mathrm{Y}}) = \{\omega_{\mathrm{X}}\}\) . If X is compact, these are simply the continuous maps from X to Y.

Let A be a commutative complex Banach algebra. Recall that \(X'(A)\) is identified with the compact space obtained from \(X(A)\) by adjoining a point at infinity (I, p.~29, corollary).

PROPOSITION 2. --- Let A and B be commutative complex Banach algebras. For every algebra homomorphism \(\pi\colon\) A \(\to\) B, the map \(\mathsf{X}'(\pi)\) is a proper partial map from \(\mathsf{X}(\mathsf{B})\) into \(\mathsf{X}(\mathsf{A})\) .

Indeed, \(\mathsf{X}'(\pi)\) is a continuous map from \(\mathsf{X}'(\mathrm{B})\) into \(\mathsf{X}'(\mathrm{A})\) (I, p.~10). The point at infinity of \(\mathsf{X}'(\mathrm{B})\) (resp. of \(\mathsf{X}'(\mathrm{A})\) ) is the zero character, and we have \(\mathsf{X}'(\pi)(0)=0\) .

PROPOSITION 3. --- a) For every proper partial map \(\varphi\) from X to Y, the map \(f \mapsto f \circ \varphi\) of \(\mathcal{C}(\mathrm{Y}')\) into \(\mathcal{C}(\mathrm{X}')\) induces an algebra morphism \(\varphi^*\) of \(\mathcal{C}_0(\mathrm{Y})\) into \(\mathcal{C}_0(\mathrm{X})\) ;

\begin{enumerate}
\def\labelenumi{\alph{enumi})}
\setcounter{enumi}{1}
\tightlist
\item
  The map \(\varphi \mapsto \varphi^{*}\) is a bijection from the set of proper partial maps from X to Y onto the set of algebra morphisms from \(\mathcal{C}_0(\mathrm{Y})\) into \(\mathcal{C}_0(\mathrm{X})\) . Its inverse bijection is the map \(\pi \mapsto \mathsf{X}'(\pi)\) .
\end{enumerate}

Let us prove a). Let \(\varphi\) be a proper partial map from X to Y, identified with a continuous map from X' to Y' such that \(\varphi(\omega_{\mathrm{X}}) = \omega_{\mathrm{Y}}\) . For \(f \in \mathcal{C}_0(\mathrm{Y})\) , one has \((f \circ \varphi)(\omega_{\mathrm{X}}) = f(\omega_{\mathrm{Y}}) = 0\) , hence the map \(\varphi^*\) is well defined. It is an algebra morphism.

Let us prove that \(\mathsf{X}'(\varphi^{*})\) is identified with \(\varphi\) . Let \(x \in X\) . For every function \(f \in \mathcal{C}_{0}(\mathrm{Y})\) , the character \(\mathsf{X}'(\varphi^{*})(\mathrm{ev}_{x})\) associates to f the complex number

\[
\left(\operatorname{ev} _ {x} \circ \varphi^ {*}\right) (f) = \operatorname{ev} _ {x} (f \circ \varphi) = f (\varphi (x)),
\]

hence \(\mathsf{X}'(\varphi^{*})(\mathrm{ev}_{x}) = \mathrm{ev}_{\varphi(x)}\) . This proves the assertion.

Conversely, let \(\pi\colon\mathcal{C}_{0}(\mathrm{Y})\to\mathcal{C}_{0}(\mathrm{X})\) be an algebra morphism. Let us prove that \(\mathsf{X}^{\prime}(\pi)^{*}=\pi\) . Let \(f\in\mathcal{C}_{0}(\mathrm{Y})\) , and denote \(g=\mathsf{X}^{\prime}(\pi)^{*}(f)\in\mathcal{C}_{0}(\mathrm{X})\) . For every \(x\in\mathsf{X}\) , one has

\[
g (x) = \left(f \circ \mathrm{X} ^ {\prime} (\pi)\right) (x) = \operatorname{ev} _ {\mathrm{X} ^ {\prime} (\pi) (x)} (f) = \left(\operatorname{ev} _ {x} \circ \pi\right) (f) = \pi (f) (x),
\]

since \(\mathsf{X}^{\prime}(\pi)\) satisfies \(\mathrm{ev}_{\mathsf{X}^{\prime}(\pi)(x)} = \mathrm{ev}_{x} \circ \pi\) . We therefore have \(g = \pi(f)\) , which allows us to conclude that \(\mathsf{X}^{\prime}(\pi)^{*} = \pi\) .

In a completely similar way, we have:

PROPOSITION 4. --- Suppose that X and Y are compact. Identify the space X (resp. the space Y) with X( \(\mathcal{C}(X)\) ) (resp. X( \(\mathcal{C}(Y)\) )) (corollary 3 of I, p.~33).

\begin{enumerate}
\def\labelenumi{\alph{enumi})}
\item
  For every continuous map \(\varphi\colon\mathrm{X}\to\mathrm{Y}\) , the map \(\varphi^{*}\colon f\mapsto f\circ\varphi\) is an algebra morphism from \(\mathcal{C}(\mathrm{Y})\) into \(\mathcal{C}(\mathrm{X})\) ;
\item
  The maps \(\varphi \mapsto \varphi^{*}\) and \(\pi \mapsto X(\pi)\) are reciprocal bijections between the set of continuous maps from \(X\) into \(Y\) and the set of algebra morphisms from \(\mathcal{C}(Y)\) into \(\mathcal{C}(X)\) .
\end{enumerate}

Remark. --- *In the language of category theory, the preceding results are interpreted as follows. Let G be the category whose objects are locally compact topological spaces and whose morphisms are proper partial maps. The functor X ↦ \(\mathcal{C}_0(X)\) is a contravariant, fully faithful functor from the category G to the category of commutative complex Banach algebras. Moreover, A ↦ X(A) is a contravariant functor from the category of commutative complex Banach algebras to the category G. If to a locally compact topological space X one associates the homeomorphism

\[
\operatorname{ev} \colon X \to X (\mathscr {C} _ {0} (X)),
\]

one obtains an isomorphism from the identity functor of the category G to the composite functor \(X \mapsto X(\mathcal{C}_{0}(X))\) .

It is not true that the composite functor \(A \mapsto \mathcal{C}_{0}(X(A))\) is isomorphic to the identity functor of the category of commutative complex Banach algebras (cf.~example 2 of I, p.~36 and exercise 2 of I, p.~155). One will nevertheless see a statement of this type for commutative stellar algebras (number 5 of I, p.~107).*

